\section{The discrete canonical profile}\label{sec:perfil}

Passing from the board to the unit square requires comparing quantities on the same scale.
Since the weights of the $n$ squares in each row sum to one, their average is $1/n$. We therefore compare
the rescaled weight $n\pi_{ij}$, rather than the original weight, with the value of $\nu_1$ at the center of
the square.

We first need to ensure that the maximizer stays away from the endpoints in the precise sense of
Section~\ref{sec:P02}: its weights lie between $c/n$ and $M/n$, and the diagonal vacancies
have a positive lower bound. Section~\ref{sec:P01} lays the groundwork for these bounds
by constructing a Slater point, that is, a feasible point satisfying the diagonal and square constraints
strictly.
We use these bounds to control the spectrum of $S_n$ in Section~\ref{sec:P03} and to expand
the local term in Section~\ref{sec:P04}. Finally, Sections~\ref{sec:P05}
and~\ref{sec:P06} establish the stated comparison with uniform error $O(1/n)$.
